The detection of gravitational waves is vital for developing our understanding of the Universe. Systems such as inspiralling binaries of black holes and neutron stars produce gravitational waves, and much of the information carried by a gravitational wave cannot be obtained via any other means. Gravitational waves interact with matter weakly, so kilometre-scale interferometers, such as the LIGO detectors, are the only instruments which have directly measured the strain induced in space-time by gravitational waves. To more accurately determine the parameters of individual sources and to refine statistical models of these systems, it is vital that the sensitivity of these interferometers is increased.
\par
To reduce the quantum shot noise of the LIGO detectors, they require a low noise, high power laser. This thesis contains experimental characterisation of the prototype for the laser that will be used during LIGO's fourth observation run. This laser generated over 100\,W of amplitude stabilised light in the HG00 mode, thus it is an important step towards reaching the design sensitivity of the LIGO detectors.
\par
Balanced homodyne detection is a key part of the upgrade from advanced LIGO to LIGO A+. To lower the quantum noise of the detectors by harnessing the quantum nature of light, it is crucial that the balanced homodyne detector has minimal loss. Mode mismatches between the interferometer and the output mode cleaners are a source of loss; therefore, active optics for mode matching between the interferometer and the output mode cleaner will be used. In this thesis, the primary sources of mismatch were explored, and the robustness of active optics to the likely types of mode mismatch was analysed.
\par
Third-generation ground-based gravitational wave detectors, such as the Einstein Telescope and LIGO Cosmic Explorer, will be far more sensitive and be able to probe deeper into the Universe than the current generation of detectors. The increase in sensitivity may be achieved with cryogenically cooled test masses, but the wavelength of light used in current gravitational wave detectors, 1\si{\micro\metre}, will not be compatible with these test masses. Instead, these test masses may work with 2\si{\micro\metre} light.
\par
High quantum efficiency photodiodes are required if the detector's quantum noise is to be minimal, so off-the-shelf extended InGaAs photodiodes that are sensitive to 2\si{\micro\metre} light where characterised in the context of the unique requirements of a gravitational wave detector. It was found that current off-the-shelf extended InGaAs will not be suitable for third-generation detectors. Significant amounts of research into the optimal conditions for manufacturing extended InGaAs photodiodes would be needed before using them in a third-generation detector is viable. 


% Photodiodes that meet the unique requirements of a gravitational wave detector have not be developed for 2\si{\micro\metre} light, and this


%To reduce the thermal noise of future detectors,, the test-masses may be cryogenically cooled. The test-mass technology a
%If future detectors, , may be operated with a laser with a wavelength of . This will require
%The detectors' quantum noise will be reduced with squeezed light, but the amount of squeezing achieved will be severely limited if balanced homodyne detection is not used.